\chapterpage{profiles}{Extract profiles from a heatmap}
\subhead{Read values along a drawn line}
Plot the example heatmap as in Chapter~\ref{heatmap}. \ui{Tools → Heatmap Profile Extractor} (\keys{Ctrl+H}) stays open while you work on the canvas.
\begin{enumerate}
\item Click \ui{Start Extraction}. On the heatmap, drag the left mouse button along 50 µm, from the first time to the last. Release to open the profile window.
\item Inspect the curve, then use its \ui{Export Data} or \ui{Save Image}. Close the result window. In the extractor, \ui{Stop} ends capture and \ui{Clear Marks} removes the line.
\end{enumerate}
\Figure{profile}{A horizontal extraction across the example. The maximum lies near the middle of the line. A manually drawn line can miss the exact centre.}{86mm}
The extractor samples 200 equally spaced points on the drawn line using the displayed heatmap grid. Its horizontal axis is distance from the start in plotted coordinate units. Along this horizontal line only time changes, so distance numerically matches elapsed ns. Along an oblique line, the calculation combines the numerical changes on both axes. Record both endpoints and their units to interpret that trace later.
\checknote{Sampling at an exact coordinate}{For an exact 50 µm entry and the original output times, use the fixed-position option on the next page. Its CSV and the drawn-line CSV have different columns and sampling.}

\nextpage{extrema}{Find extrema on a defined section}
\begin{enumerate}
\item With the heatmap active, open \ui{Tools → Extreme Finder Tool} (\keys{Ctrl+E}). Choose \ui{Fixed Position}, enter \code{50}, then click \ui{Add Profile}. Comma-separated positions add several rows.
\item Keep the full range, enable \ui{Find Maxima} and click \ui{Extract \& Find Extremes}. Open the profile-curve tab. For the example, the maximum is 8 eV at 2 ns.
\end{enumerate}
\Figure{extrema}{The marked maximum belongs to the time history at 50 µm. The profile table identifies the requested position and any search ranges.}{86mm}
\subhead{Position cuts and time cuts}
A fixed-position cut interpolates spatially onto that position at each recorded time. For a fixed-time cut, the viewer takes the nearest stored time within the file's time span and uses that record's coordinates. Enter an actual output time when the distinction matters.
Edit a profile's range cell to restrict the search, for example \code{1-3} ns for the 50 µm cut. The finder reports a maximum or minimum per specified interval. Its derivative view uses numerical gradients. Near gaps, broad peaks or boundaries, compare against the plotted samples. Save the active image or export the profile table to retain the values as well as the marker.
\checknote{Table columns}{The extrema CSV uses \code{X} and one column per requested section. Export position cuts and time cuts separately so that every value in the \code{X} column has the same physical unit.}

\chapterpage{annotations}{Annotate and inspect a plot}
\subhead{Label a feature for a report}
With the three spatial profiles visible, press \keys{Ctrl+M}. Choose \ui{Arrow}, set colour and line width, and drag from the label side towards the 2 ns peak. Choose \ui{Text}, click beside the arrow and enter \code{2 ns}. Use the annotation list to select and delete an individual mark.
\Figure{annotation}{An arrow and a text label identify the central peak on the plotted curves. The saved image can carry the mark without changing any of the underlying values.}{72mm}
The annotation palette also offers a straight line, rectangle and ellipse, style, transparency and text controls. Use \ui{Export Annotated Image} while the marks are visible. Replotting rebuilds the axes, so export the annotated figure before changing the plot. The tab's numerical CSV contains the underlying values.
\subhead{Inspect a detected trajectory}
With a suitable dataset loaded, activate the heatmap and choose \ui{Tools → Shock Wave Detector} (\keys{Ctrl+D}). The dialog can show the field, profiles and a detected trajectory, with controls for inspecting results, saving the heatmap and exporting a trajectory. Check the indicated position against the source profile at the same time. Keep the field, axes, units and detector configuration with any exported measurement. The practice case contains a smooth temperature peak and demonstrates the plotting tools.
\checknote{Mouse interaction}{Close the annotation or extraction tool to release the plot for ordinary drag-to-zoom navigation.}

\chapterpage{export}{Save figures and numerical data}
Use \keys{Ctrl+S} on the active plotting tab to save the current figure. Choose PNG for a raster image or PDF for a figure to place in a document. Type the extension explicitly. The main plotting tabs use a default save resolution of 300 dpi for raster output.
\Table{Command}{Contents of the output}{
Spatial tab → Export Data & One coordinate column. Each enabled, selected moment adds a field column. The time in its header is in ns.\\
Temporal tab → Export Data & One time column in the selected unit. Global mode adds one scalar column. Cell mode adds one column per enabled, selected cell.\\
Heatmap tab → Export Data & For an ordinary spatial field, one coordinate column followed by one per stored time. The exporter resamples onto 500 coordinate rows. A boundary spectrum follows a separate frequency-group export path.\\
Profile window → Export Data & Two columns, \code{Distance} and the variable key, with 200 samples along the drawn line. Record the axis orientation, endpoints and units alongside it.\\
Extreme finder → Export Data & Shared \code{X} column and one value column per section. Blank fields may occur for positions absent in a section.\\}
\subhead{Export menu commands}
\ui{Export → Export as CSV} creates a separate file for each of the first five available spatial variables, or fewer when the file has fewer variables. Columns represent timesteps, rows follow array order, and the data come from the parser in stored units. Use the source file to match array rows to physical coordinates. Give the base name a lowercase \code{.csv} suffix, for example \code{run.csv}, to obtain names such as \code{run_T.csv}.
\ui{Export → Export as ASCII} uses the legacy writer. Enter a base name such as \code{run.txt} to create \code{run_SCout.txt} and \code{run_TSout.txt}. The writer expects its classic field set. For a reduced dataset such as the tutorial, use the plotting tab's CSV when a legacy field is missing.
\checknote{Before handing off data}{Save into a writable folder under a distinct basename. Record the original filename, variable, coordinate family, units, selected times or cells, and any interpolation or axis limits.}

\nextpage{checkcsv}{Check exported values}
Keep an unchanged copy of \code{quickstart.zip}. Create a new output directory. From the spatial, cell and heatmap tabs save \code{profiles.csv}, \code{cell.csv} and \code{map.csv}. Open each as a comma-separated UTF-8 table. The row counts below exclude the header.
\Table{File}{Checks for the practice case}{
\code{profiles.csv} & 21 rows and four columns for moments \#0, \#4, \#8. At coordinate 50 µm the 2 ns column is 8.000000e+00 eV.\\
\code{cell.csv} & 9 rows and two columns for cell \#10. The time column reads 2.000000e+00 ns on the fifth data row and the matching value is 8.000000e+00 eV.\\
\code{map.csv} & 500 coordinate rows and ten columns. The first column spans 0 to 100 µm. Since 500 is even, the two middle rows lie on either side of 50 µm.\\
Global temporal CSV & 9 rows and two columns, \code{Time (ns)} and \code{ENIE}. At 2 ns the scalar is 14 in stored units.\\}
\subhead{Example header from the spatial tab}
\begin{verbatim}
XC Centroid (um),t=0.0000ns T (eV),t=2.0000ns T (eV),
t=4.0000ns T (eV)
\end{verbatim}
The header is wrapped here for readability. A CSV reader sees four fields on one line. Import the first row as column names and retain scientific notation when copying values.
\subhead{When the mesh moves}
In version 2.3.1 the spatial CSV uses the first stored coordinate vector for its single position column. The ordinary heatmap CSV likewise resamples each field record using that reference grid. The on-screen Eulerian heatmap uses each record's coordinates. For quantitative work on a moving grid, pair the original field with \code{XC} or \code{X} at the same time from the source output. For a fixed physical position across time, use the fixed-position section in the extreme finder and retain its axis units.
\checknote{Why a value can differ slightly}{The plotted fixed-time profile has original cell samples. A heatmap grid point and a point on a drawn line may fall between two centres. Compare at matching coordinates, times and value units.}

\chapterpage{recorder}{Record and replay a plotting session}
Use \code{quickstart.out} for this exercise. The generated Python script embeds the classic binary parser. The GUI can also read decimal cases and formatted inputs. For those inputs, save figures and numerical tables directly from the GUI.
\begin{enumerate}
\item Open \code{quickstart.out}. Press \keys{Ctrl+R} and click \ui{Start Recording}. The recorder includes the active file as an operation.
\item Return to \ui{Spatial Plot}, choose \ui{T}, \ui{XC}, \ui{um}, \ui{eV} and moment \ui{\#4}. Click \ui{Plot}. Reopen the recorder if needed, click \ui{Stop Recording}, then \ui{Generate Script}.
\item Save the script as \code{example_replay.py} in a writable folder with \code{quickstart.out}. Change to that directory in a terminal with Python, NumPy, SciPy and Matplotlib installed.
\end{enumerate}
\begin{verbatim}
python example_replay.py
\end{verbatim}
\Figure{recorder}{After recording one plot on an already loaded file, the counter shows two operations. Generate Script is now available.}{54mm}
The example produces \code{example_replay_fig1_spatial_plot.png} in the working directory. The script first tries the recorded absolute input path, then a file of the same name beside the script. Open the PNG and verify a single 2 ns profile with a maximum near 50 µm. Inspect the exported CSV separately if you need the exact cell value.
\checknote{Starting a new recording}{A new recording clears the prior operation list. Generate and name the previous script first. A plotting replay captures supported operations. Save interactive markups separately as annotated images.}

\chapterpage{trouble}{Troubleshoot a session}
\Table{Symptom}{What to check}{
Double-click shows package contents & Choose the system software installer in \ui{Open With}, or run the local-file APT or DNF command in Chapter~\ref{install}. Review any dependency error printed by the package manager.\\
The launcher exits or no window appears & Start \code{multi1dvisual} from a terminal in the desktop session. Check the first Qt or missing-package message. A remote shell needs access to a graphical display.\\
Chinese characters look empty & Install a Chinese font through the distribution's font manager and restart the viewer. On Debian or Ubuntu the recommended font package is \code{fonts-noto-cjk}.\\
A decimal case fails to open & Extract the archive first. Check the basename and letter case of both required tables, tab delimiters, the \code{TIME} header, and the number and ordering of rows per time.\\
A curve is absent or nearly flat & Check the active file, variable, selection and the exclusion state. Enable automatic axis ranges, clear log scaling, click \ui{Plot}, then check the source column and units.\\
A heatmap is blank or saturated & Choose linear colour, recalculate the range with \ui{Auto}, and check axis limits. Inspect an original spatial record for finite values. Gaps can occur outside the recorded spatial support.\\
A tool asks for a heatmap & Activate the heatmap tab, select a supported field, plot it and then open the tool. Finish a previous annotation or line extraction if dragging has the wrong effect.\\
A PDF manual fails to open & Check the installed documentation directory and its PDF files. Open one directly with a local PDF reader and check your desktop's default PDF application.\\}
\subhead{Find a diagnostic log}
In your home directory, open the hidden \textsf{.MULTI1D可视化工具} folder and its \code{logs} subfolder. Recent runs use filenames beginning with \code{run_}. The logger retains the three newest log files. For a reproducible report, include the application version, operating system, input format, selected view and units, the sequence that failed, and the relevant lines from the latest log.

\chapterpage{reference}{Shortcuts and file reference}
\Table{Shortcut}{Action}{
\keys{Ctrl+O} / \keys{Ctrl+S} / \keys{Ctrl+Q} & Open an input file / save the active figure / quit.\\
\keys{Ctrl+1} / \keys{Ctrl+2} / \keys{Ctrl+3} & Switch to spatial, temporal or heatmap view.\\
\keys{Ctrl+M} & Open the annotation palette for the active canvas.\\
\keys{Ctrl+H} / \keys{Ctrl+E} & Open the drawn-line extractor / extreme finder from the heatmap.\\
\keys{Ctrl+D} / \keys{Ctrl+R} & Open the shock detector from the heatmap / open the code recorder.\\
\keys{A} / \keys{H} & With the canvas active, restore the full view / step back one zoom level.\\}
\subhead{Common field names}
\Table{Stored key}{What the label means}{
\code{TIME}, \code{XC}, \code{X} & Output time in seconds, cell-centre position in cm and interface position in cm before display conversion.\\
\code{T}, \code{TI}, \code{TR} & Electron, ion and radiation temperatures respectively. The example provides \code{T}.\\
\code{R}, \code{P}, \code{PT} & Mass density, pressure and total pressure respectively. Units on the plotted axes depend on the field and display choice.\\
\code{ENIE}, \code{ENII}, \code{ENKI} & Global electron internal energy, ion internal energy and kinetic energy. Check the producer's units when a field label omits them.\\}
\subhead{Preferences and project files}
The application uses \code{config.json} and \code{language_config.json} in the per-user directory described on the previous page. The project README links both manuals and the example. The archive contains \code{quickstart.case}, \code{quickstart.center.dat}, \code{quickstart.scalars.dat} and \code{quickstart.out}. To see exactly how the synthetic numbers were made, open \code{docs/examples/make_quickstart.py} in the source tree.
\checknote{Project}{Software by Bo Xu. The repository contains the source, package build instructions and editable sources for these manuals. \repo}
